% SPDX-License-Identifier: Apache-2.0
\section{The Rule, and What It Leaves}
\label{sec:e-rule}

\begin{rulebox}[The reduction]
Delete every construct Rust already provides. Rename every construct whose
name collides with a Rust concept. What remains is the library.
\end{rulebox}

The renaming clause does one job in practice. A hardware \term{module} and
a Rust \code{mod} are different things with one name, so the hardware one
becomes a \term{unit} and \code{mod} keeps its meaning as a namespace.

\subsection{The vocabulary that remains}

Two Rust modules hold everything the language still needs.

\begin{lstlisting}
crate::types::Transaction   // data that moves between units
crate::types::Tag           // a domain and its policy

crate::comp::Module         // a unit
crate::comp::Bus            // a bundle of channels
crate::comp::Chan           // one channel of a bus
crate::comp::Signal         // one wire
crate::comp::Port           // an attachment to a bus, in a role
crate::comp::Role           // the type-level tag for a modport
\end{lstlisting}

Everything else is a plain Rust struct, trait or function.
Table~\ref{tab:reduction} states the mapping.

\begin{table}[H]
\caption{What each construct becomes. Seventeen of twenty are deleted
outright.}
\label{tab:reduction}
\centering
\footnotesize
\begin{tabular}{@{}ll@{}}
\toprule
Was & Is now \\
\midrule
\code{transaction T}   & \code{struct T} plus \code{impl Transaction} \\
\code{module M}        & \code{struct M}, a unit, plus \code{impl Module} \\
\code{bus B}           & \code{struct B} plus \code{impl Bus} \\
\code{interface I}     & \code{struct I} of \code{Signal} fields \\
\code{modport R}       & \code{trait R}, implemented for the interface \\
\code{channel c: T}    & a field of type \code{Chan<T>} \\
\code{port p}          & a field of type \code{Port<B, R>} \\
\code{flow f}          & \code{fn f} \\
\code{seq s}           & \code{async fn} \\
\code{pipeline stage pipe} & \code{async fn} and \code{.await} \\
\code{tag T}           & \code{struct T} plus \code{impl Tag}, as a parameter \\
\code{config bind}     & a \code{type} alias \\
\code{domain D}        & a type parameter \\
\code{assert assume cover} & Rust's own macros \\
\bottomrule
\end{tabular}
\end{table}

\subsection{Transactions, interfaces and modports}

A transaction is a struct that implements a marker trait. Nothing else
about it is special.

\begin{lstlisting}
pub struct MacRequest { pub addr: u32, pub count: u8 }
impl Transaction for MacRequest {}
\end{lstlisting}

An interface splits into two Rust items, and that split fixes a defect the
earlier analysis reported~\cite{unification}. LHDL's \code{interface} did
two jobs at once, a bundle of wires and a socket an implementation fills,
which is why its grammar and its examples never agreed. Here the bundle is
a struct and the role is a trait.

\begin{lstlisting}
// The wires. No direction on any of them.
pub struct Wishbone {
    pub adr: Signal<u32>,
    pub ack: Signal<bool>,
}

// A role. Direction lives here, not in the struct.
pub trait Initiator {
    fn drive_adr(&mut self, v: u32);
    fn read_ack(&self) -> bool;
}

impl Initiator for Wishbone { /* ... */ }
\end{lstlisting}

A unit may then be generic over the role it attaches in rather than over a
concrete bus, which is what late binding needs and
Section~\ref{sec:e-binding} develops.
